% Management in Global Markets – Midterm All-in-One (Concepts, Cases, Exam Tips)
% Text-to-speech friendly version (simple sections, no columns or math symbols)

\documentclass[10pt,a4paper]{article}
\usepackage[utf8]{inputenc}
\usepackage[T1]{fontenc}
\usepackage[margin=2cm]{geometry}
\usepackage{enumitem}
\setlist{nosep,leftmargin=1.5em}
\raggedright

\begin{document}

\section*{Management in Global Markets Midterm All-in-One}

This sheet is built from your course notes and the week six midterm guide. It is meant to be listened to, not printed pretty.

\subsection*{How the exam works}
\begin{itemize}
\item Midterm: closed notes, open questions, sometimes based on one slide with a graph or drawing. You answer in short, clear lines by hand.
\item The professor does not want long essays; he wants you to write one or two strong lines per point.
\item Your job in each answer: name the concept, define it in a short sentence, give one example from class or slides, then explain why it matters for international business.
\end{itemize}

\section*{1. Big ideas of the course}
\begin{itemize}
\item International business means operating across borders. It involves products, services, capital, and people.
\item The central question is: why go abroad at all? Typical answers are larger markets, access to resources, lower costs, and learning from other contexts.
\item In reality, most companies fail in internationalisation because of wrong location, weak competitive advantage, or poor execution.
\item Key sentence from the professor: “Location is the masterpiece.” The choice of place is often more important than skills, timing, or effort.
\end{itemize}

\section*{2. Globalisation, slowbalisation, and economic complexity}
\begin{itemize}
\item Globalisation means increasing flows of goods, services, capital, and sometimes people across borders; economies become highly interconnected.
\item Slowbalisation describes a slowdown in these flows, with more protectionism and localisation.
\item Winners from globalisation include China, India, and Indonesia; the United States is the biggest importer and China a major exporter.
\item Economic complexity means moving from simple raw products to complex products that require skills and know‑how; countries with high economic complexity have more advanced, diversified exports.
\item The course stresses that rich countries tend to have high economic complexity and that complexity comes from skills, training, and technology, not from natural resources alone.
\end{itemize}

\section*{3. Supply chains and products}
\begin{itemize}
\item A supply chain is the sequence of steps and locations that parts go through before becoming a final product; rarely is a product “made” in one country only.
\item Example from class: shoes from Vietnam, shirts from India, watches from China, or aircraft assembled in Chicago with parts from many countries.
\item Japan is associated with quality and higher price; China with lower cost and high volume; quality ranges widely.
\item International supply chains create risks: cultural, legal, logistical, financial, and communication; you cannot escape language and culture.
\item Market attractiveness involves growth in GDP, population, and an assessment of political and economic risk.
\end{itemize}

\section*{4. Culture and Hofstede dimensions}
\begin{itemize}
\item Culture is defined as shared beliefs, values, and norms that distinguish one group from another.
\item Hofstede’s dimensions include power distance (acceptance of hierarchy), individualism versus collectivism, uncertainty avoidance, and long‑term versus short‑term orientation.
\item Rich countries tend to be more individualistic; poorer countries more collectivist; but there are many nuances.
\item The key idea is that you cannot escape culture; you must adapt your behaviour and expectations when entering new markets.
\item Culture shock typically has stages: initial enthusiasm, then difficulty, then adaptation.
\end{itemize}

\section*{5. Four phases of internationalisation and location risk}
\begin{itemize}
\item Phase one, conceive: you plan, research, and evaluate markets. This can take two or three years.
\item Phase two, execute: you invest and open plants or operations; this is the stage of highest investment and risk.
\item Phase three, establish: you try to stabilise operations, reach profitability, and make the project sustainable; the professor calls this the most dangerous phase because many projects fail here.
\item Phase four, expand: you grow once the model works.
\item Artificial risks come from laws and rules and can be managed; natural risks come from language and culture and must be adapted to.
\item Many failures come from choosing the wrong location in phase two or three, even when the idea and product look strong.
\end{itemize}

\section*{6. Strategy levels and competitive advantage}
\begin{itemize}
\item Corporate strategy answers: in which businesses will we compete? For example, cars versus trucks.
\item Competitive strategy answers: how will we compete in each business? For example, price versus quality versus niche.
\item Functional strategy refers to decisions in HR, marketing, finance, operations, and research and development.
\item Small and Medium Enterprises, or SMEs, have fewer than two hundred and fifty employees and less than fifty million in turnover; they are the majority of firms but few successfully internationalise.
\item Resources are what the company has; capabilities are what the company does with those resources.
\item Competitive advantage means being clearly better than rivals; absolute competitive advantage, where nobody can match you, is rare.
\item Before going abroad, companies must check if their advantage is strong enough; going out as “just another similar firm” in a mature market is extremely risky.
\end{itemize}

\section*{7. Chile case and technology transfer}
\begin{itemize}
\item In the Chile case from the course, a French company brought capital, people, and technology to a partnership in Chile.
\item The project failed: the company lost money, jobs, and technology; the partners used the technology independently afterwards.
\item There was weak protection of intellectual property and poor communication and trust.
\item Lesson: be very careful with core technology when entering new markets; choose partners wisely; use contracts and protections; understand culture and build trust.
\end{itemize}

\section*{8. Export, international operations, and FDI}
\begin{itemize}
\item Export means the product moves to other countries but production stays at home; risk and commitment are lower.
\item International operations mean the company moves plants or key functions abroad; this is a deeper commitment.
\item Foreign Direct Investment means a long‑term commitment of capital and people in a foreign country; it increases complexity in human resource management and risk.
\item A simple progression is: domestic first, then export, then international operations with FDI.
\end{itemize}

\section*{9. International human resources}
\begin{itemize}
\item International HRM means managing people in several countries and cultures and coordinating policies across borders.
\item Questions you must answer include: which functions move, how many people do you need, do you have them already or must you hire, and how will you support families.
\item Often younger and more flexible workers go abroad first; for sensitive positions companies prefer internal people they already trust.
\item Motivation is not just about salary; it also involves ambition, recognition, and quality of life.
\item Turnover among young workers is high; you have to know their market value and negotiate.
\item Trust is central: giving “keys to the kingdom” to outsiders is risky, which is why many firms prefer to use people they know.
\end{itemize}

\section*{10. Cash generation}
\begin{itemize}
\item Cash ultimately comes from customers; if you do not sell, there is no cash.
\item The two classic options for improving profit are: reduce expenses, or push sales.
\item Repeat business is extremely important; communication and delivering on promises are key to repeat sales.
\end{itemize}

\section*{11. Exam checklist}
\begin{itemize}
\item Can you explain in your own words: globalisation, slowbalisation, and economic complexity?
\item Can you describe a simple supply chain and explain why products are “from many countries”?
\item Can you name at least two Hofstede dimensions and say how they change business behaviour?
\item Can you describe the four phases of internationalisation and say which is the most dangerous and why?
\item Can you explain the difference between export, international operations, and foreign direct investment?
\item Can you describe the Chile case and its lesson about technology and partners?
\item Can you explain why “Location is the masterpiece” and give an example?
\end{itemize}

\end{document}

